%document class
\documentclass[10pt,twocolumn]{book}
%%%% Page Info + Commands %%%%%{

%packages
\usepackage{pgfplots}
\pgfplotsset{compat=1.18}
\usepackage{amsmath}
\usepackage{amsfonts}
\usepackage{charter}
\usepackage{amsthm,letltxmacro}
\usepackage{amssymb}
\usepackage{fancyhdr}
\usepackage{float}
\usepackage[most]{tcolorbox}
\usepackage{enumerate}
\usepackage{enumitem}
\usepackage[dvipsnames]{xcolor}
\usepackage{changepage}
\usepackage{mdframed}
\usepackage{graphicx}

% SMALL PAGE COMMAND
\newcommand{\smallpage}{  \setlength\oddsidemargin{.25in}
            \setlength\textwidth{6in}
            \setlength\topmargin{0in}
            \setlength\textheight{9in}}


% Commands for sets of numbers
\newcommand{\Z}{\mathbb{Z}}\newcommand{\R}{\mathbb{R}}\newcommand{\N}{\mathbb{N}}\newcommand{\Q}{\mathbb{Q}}\newcommand{\C}{\mathbb{C}}
\newcommand{\real}[1]{\text{Re}\left(#1\right)}
\newcommand{\im}[1]{\text{Im}\left(#1\right)}

% gordie spacing and boxing commands
\newcommand{\gspc}{\vspace{0.13cm}} % 0.5 space
\newcommand{\gline}[0]{\gspc\\} % 1.5 space
\newcommand{\gbline}{\gspc\gspc\\} % double space
\newcommand{\gbox}[1]{\fbox{\parbox{\linewidth}{#1}}} % multiline box command
\newcommand{\gmod}[1]{\,(\text{mod}\,#1)}
\newcommand{\gbar}[1]{\overline{#1}}
\newcommand{\cpr}[1]{\left(#1\right)}
\newcommand{\ccbr}[1]{\left\{#1\right\}}
\renewcommand{\Re}[1]{\text{Re}\left(#1\right)}
\renewcommand{\Im}[1]{\text{Im}\left(#1\right)}
\newcommand{\cis}{\,\mathrm{cis}\,}
\newcommand{\Arg}{\mathrm{Arg}\,}
\newcommand{\ahline}{\\[3pt]\hline\\[-8pt]}


\newcommand{\gimage}[2]{\vspace{-0.25cm}
\begin{figure}[H]
    \centering
    \includegraphics[width=#2]{#1}
\end{figure}\vspace{-0.4cm}
}

\newcommand{\centerit}[1]{{\begin{center} #1 \end{center}}}
\newcommand{\ul}[1]{\underline{#1}}

%%%%%%%% command for graphics %%%%%%%%%%%%%

%}

% COLOR BOX
\newmdenv[backgroundcolor=gray!8,
  linewidth=0pt, innerleftmargin=0pt, innerrightmargin=0pt,
  skipabove=\baselineskip, skipbelow=\baselineskip
]{coloredadjust}

% PROOF
\LetLtxMacro\oldproof\proof
\let\endoldproof\endproof
\renewenvironment{proof}[1][\proofname]
  {\vspace{0.2cm}\begin{adjustwidth}{2em}{2em}
   \oldproof[#1]}
  {\endoldproof
\end{adjustwidth}}

% PROBLEM
\newenvironment{problem}[1]
  {\vspace{-0.1cm}\begin{coloredadjust}\begin{adjustwidth}{2em}{2em}#1}
  {\endoldproof
\end{adjustwidth}\end{coloredadjust}\gspc}

%% DEFINITION
\newtcbtheorem[]{definition}{Definition}{enhanced,
  colback=gray!2, colframe=gray!50!black, coltitle=gray!30!black,
  fonttitle=\bfseries, breakable,
  boxrule=0pt, toprule=2pt, bottomrule=1pt, leftrule=1pt, rightrule=1pt, arc=1pt, outer arc=1pt,
  attach title to upper, title={#1},
}{dfn}

%% CRITERION
\newtcbtheorem[]{criterion}{Criterion}{enhanced,
  colback=magenta!3, colframe=magenta!60!black, coltitle=magenta!20!black,
  fonttitle=\bfseries,
  boxrule=4pt, toprule=2pt, bottomrule=1pt, leftrule=1pt, rightrule=1pt, arc=1pt, outer arc=1pt,
  attach title to upper, title={#1},
}{ctn}

%% COROLLARY
\newtcbtheorem[]{corollary}{Corollary}{enhanced,
  colback=cyan!2, colframe=cyan!60!black, coltitle=cyan!30!black,
  fonttitle=\bfseries,
  boxrule=4pt, toprule=1pt, bottomrule=1pt, leftrule=1pt, rightrule=1pt, sharp corners,
  attach title to upper, title={#1},
}{ctn}

%% THEOREM
\newtcbtheorem[]{theorem}{Theorem}{enhanced,
  colback=blue!2, colframe=blue!60!cyan!30!gray, coltitle=blue!40!cyan!50!black,
  fonttitle=\bfseries,
  boxrule=4pt, toprule=4pt, bottomrule=2pt, leftrule=1.5pt, rightrule=1.5pt,
  sharp corners,
  breakable,
  attach title to upper, title={#1},
}{thm}

\newenvironment{example}{
  \begin{tcolorbox}[colback=green!2, colframe=green!40!black, coltitle=green!20!black, sharp corners]
    \textit{Example. }
}{\end{tcolorbox}}


%%%% Page Setup %%%%%%%
\smallpage\sloppy
\pagestyle{fancy}
\lhead{\large{\textbf{Complex Analysis Notes}}}
%\rfoot{\thepage/\pageref{LastPage} }
\setlength{\headheight}{10pt}

\begin{document}
\small

\subsection*{1.2 Point Representation of Complex Numbers}
There is a natural bijection between $\C$ and the $xy$-coordinate system $\R^2$:
\begin{align*}
  a + bi \leftrightarrow (a,b)
\end{align*}
\begin{theorem}{Properties of Complex Nums}{}
  \begin{enumerate}
    \item $\gbar{ x \pm w} = \bar z \pm \bar w$
    \item $\gbar{zw} = \bar z \cdot \bar w$
    \item $\gbar{\cpr{\frac{z}w}} = \frac{\bar z}{\bar w}$
    \item $\Re{z} = \frac{z + \bar z}{z}$
    \item $\Im{z}= \frac{z - \bar z}{zi}$
  \end{enumerate}
\end{theorem}
\begin{definition}{Modulus}{}\gline{}
  The \ul{modulus} of $z = a+bi$ is:
  \begin{align*}
    |z| = \sqrt{a^2 + b^2}
  \end{align*}
  We can imagine this as the distance from the origin:
  \gimage{image1.png}{3cm}
\end{definition}

\begin{theorem}{Properties of the Modulus}{}
  \begin{enumerate}
    \item $|z| = |\bar z|$
    \item $|z| = z$ iff. $z$ is a nonnegative real number
    \item $z\bar z = |z|^2$
    \item $|z+w| \le |z| + |w|$
    \item $|z-w| \ge ||z|-|w||$
    \item $|zw| = |z|\cdot|w|$
    \item $\left|\frac{z}{w}\right|=\frac{|z|}{|w|}$
  \end{enumerate}
\end{theorem}
\subsection*{1.3 Vector and Polar Forms}
We saw last time that complex addition behaves like vector addition.\gline{}
Now, consider \ul{Polar Form}:
\gimage{image3.png}{7cm}
More formalized, for \textit{any $z\in\C$ of the form $a + bi$}, we can write:
\begin{align*}
  z = &= a + bi\\
  &= r(\cos\theta +i\sin\theta)\\
  &= r\cis\theta
\end{align*}
In this case, we write $r = \sqrt{a^2 + b^2}=|z|$. Note that $\theta$ is defined as counter-clockwise from the positive real axis. \gline{}
Here's some quick conversion equations:
\begin{align*}
  \mathrm{Polar}&\to\mathrm{Cartesian}\\
  x &= r\cos\theta\\
  y&= r\sin\theta\\
  \mathrm{Cartesian}&\to\mathrm{Polar}\\
  r &= \sqrt{x^2 + y^2}\\
  \theta &= \tan^{-1}\cpr{\frac{y}x}^*
\end{align*}
\textit{*Note: If $\theta$ is actually in the 2nd or 3rd quadrant, you'll need to add $\pi$ to $\tan^{-1}\cpr{\frac{y}x}$ to get a valid angle for $\theta$}\gline{}
\begin{definition}{Argument of $z$}{}\gline{}
  The \ul{argument of $z$}, denoted $\arg(z)$, is the set ofall positive values of $\theta$ in the polar form of $z$. Note that $\arg(z)$ is undefined for $z = 0$.  
  \begin{example}
    Let $z = -3 -3i$. Then, $\arg(z)=\ccbr{\frac{5\pi}4 + 2\pi k\mid k\in \Z}$.
  \end{example}
\end{definition}
\begin{definition}{Principal Argument of Z}{}\gline{}
  The \ul{principle oargument of $z$}, denoted $\Arg{z}$, is the element of $\arg(z)$ whose value is in the interval $(-\pi, \pi]$.
  \begin{corollary}{Branch of $\arg(z)$}{}\gline{}
    The notation $\arg_\tau(z)$ will denote the element of $\arg(z)$ whose value lies in the interval $(\tau, \tau + 2\pi]$. This is called a \ul{branch} of $\arg(z)$. 
  \end{corollary}
\end{definition}

\subsubsection*{Complex Multiplcation}
Consider the following algebra:
\begin{align*} 
  \mathrm{If }z &= r(\cos\theta+i\sin\theta)\mathrm{ and }w=s(\cos\phi+i\sin\phi)\\
  zw&= r (\cos\theta+i\sin\theta)s(\cos\theta+i\sin\phi)\\
  &= rs((\cos\theta\cos\phi-\sin\theta\sin\phi))\\&\quad+i((\sin\theta\cos\phi + \cos\theta\sin\phi))\\
  &= rs(\cos(\theta+\phi)+i\sin(\theta+\phi))
\end{align*}
\subsection*{1.4 The Complex Exponential}
First, we begin with Euler's formula and how we can formulate powers of imaginary numbers. 
\begin{definition}{Euler's Formula}{}\gline{}
  Given $y\in\R$, we define $e^{iy}$ to be:
  \begin{align*}
    e^{iy}=\cos y + i \sin y \text{ (Euler's Formula)}
  \end{align*}
  More generally, for $z = x + iy$, we define:
  \begin{align*}
    e^{z}&= e^x(\cos y + i\sin y)
  \end{align*}
\end{definition}
Note that this function agrees with the formula for real-valued inputs, because if $iy = 0$, then $e^z = e^x$. If you remember Taylor or Maclauren series, you may note that the expansion of $e^x$ is deeply related to sines and cosines, and this is that relationship embodied. \gline{}
This is also a much shorter way to use polar coordinates. \gline{}From now on, we'll write the polar form of any given $z$ as:
\begin{align*}
  x = re^{i\theta}, \text{ where }r = |z| \text{ and }\theta\in\arg(z)
\end{align*}
\begin{example}
  Let $z = \sqrt 3 + i$. Then:
  \begin{align*}
    e^z = e^{\sqrt 3 + 1} &= e^{\sqrt 3}\cdot e^i\\
    &\approx 5.652\cdot e^{i}
  \end{align*}
  So, $e^{\sqrt 3 + i}$ has $r = e^{\sqrt 3}$ and $\theta = 1$. 
\end{example}
We can formulate some general properties for multiplication as well. Consider this formulation:
\begin{align*}
  z &= re^{i\theta} \text{ and } w= se^{i\phi}, \text{ then}\\
  zw&= (re^{i\theta}(se^{i\phi}))\\
  &= rse^{i\theta + \phi}\ahline
  \frac{z}w &= \frac{r}se^{i(\theta - \phi)}\\
  z^n&= r^ne^{in\theta}\text{, for }n\in\Z\\
  \frac{1}z&= \frac{1}re^{-i\theta}
\end{align*}
\subsubsection*{1.4 Exercises - Complex Exponentials}
\begin{enumerate}
  \item Write the following complex numbers in polar form $re^{i\theta}$.
  \begin{enumerate}
    \item $e^{e^i}$
    \begin{problem}
      \begin{align*}
        e^{\cos(1) + i\sin(1)}&= e^{\cos 1}e^{i\sin 1}\\
        &= e^{\cos 1}(\cos\sin 1 \\&\quad\quad+ i \sin\sin 1)
      \end{align*}
    \end{problem}
    \item $(2 - 3i)^4$
    \begin{problem}
      \begin{align*}
        |2 - 3i| = \sqrt{13}\\
        \arctan\cpr{\frac{-3}2} = -0.98\ahline
        (\sqrt{13}e^{-0.98})^4 = 169e^{i(-3.92)}
      \end{align*}
    \end{problem}
    \item $\frac{1+\sqrt 3i}{2-2i}$
    \begin{problem}
      \begin{align*}
      \frac{1+\sqrt 3i}{2-2i}&= \frac{2e^{i\frac{\pi}3}}{2\sqrt 2 e^{i\cpr{\frac{-\pi}4}}}\\
      &= \frac{\sqrt 2}2 e^{i\frac{\pi}2}
    \end{align*}
    \end{problem}
  \end{enumerate}
\end{enumerate}
There are some more useful facts that we can utilize to easily take advantage of Euler's formula.
\begin{theorem}{Properties of Euler's Formula}{}
  \begin{enumerate}
    \item $|e^{i\theta}|=1, \forall \theta \in \R$
    \item $|re^{i\theta}|=r$
    \item $e^0=e^{i0}=1,\quad\quad e^{i(2\pi k)}=1$
    \item $e^{i\frac{\pi}2}= i$
    \item $\cos\theta = \Re{e^{i\theta}}=\frac{e^{i\theta + e^{-i\theta}}}2$
    \item $\sin \theta = \Im{e^{i\theta}}$
    \item $e^{i\pi}=-1\quad\text{ or }\quad e^{i\pi}+1=0$
  \end{enumerate}
\end{theorem}
Now, there's also one more formula that we should be aware of that comes from the textbook as well. It's kind of a relation to Euler's equation.
\begin{theorem}{DeMoirre's Formula}{}\gline{}
  $\forall\theta\in \R, \text{ and } \forall n \in \N$, 
  \begin{align*}
    \cos(n\theta)+i\sin(n\theta)
  \end{align*}
\end{theorem}
\subsubsection*{\ul{$e^{z}$ as a mapping}}
We can imagine $e^z$ as a mapping from imaginary to real. For example,
\gimage{image4.png}{6cm}









\end{document}